\section{Spekkens' diagnosis and a Leibnizian layer principle}\label{sec:spekkens}

\subsection{The category mistake: inference versus causation}

A recurring source of confusion in quantum foundations is a silent slide between two kinds of statements:
(i)~\emph{causal} statements about what the world does, and
(ii)~\emph{inferential} statements about what an agent (or an interface) can stably record, compute, or predict.
Quantum theory, as normally presented, uses a single mathematical object---the quantum state---to play both roles at once. This conflation is not merely a matter of interpretation; it shapes what counts as a paradox.

Spekkens' central diagnosis is that physicists are making a category mistake: they treat an inferential/completion object (the quantum state) as if it were an element of physical ontology.\cite{Spekkens2007Toy} Once this conflation occurs, standard arguments pressure theorists toward surplus structure---distinctions that are posited to exist in reality but are invisible to any admissible record at the descriptive layer under consideration. The result is familiar: discontinuous collapse, measurement-induced ``disturbance'' reified as a physical shock, and nonlocal influences that cannot be used to signal and hence cannot be isolated as operationally distinguishable causal mechanisms at the relevant layer.

SBT's starting point is to refuse this conflation. SBT separates substrate dynamics (causal evolution) from packaging (the process that turns some distinctions into stable record-level objects). To motivate this separation as a methodological principle rather than a mere preference, we adopt a Leibnizian stance.

\subsection{OI--EI as a methodological constraint}

\begin{principle}[Ontological identity of empirical indiscernibles (OI--EI)]
Fix a descriptive layer (a choice of admissible lenses/records/experiments). If two putative scenarios are empirically indiscernible under \emph{all} admissible experiments at that layer, then those two scenarios must be identified as the same object of description at that layer.
\end{principle}

This identification is layer-relative (analogous to quotienting gauge redundancy): OI--EI does not deny microstructure; it denies only that empirically silent distinctions are objects of the chosen layer.
The principle is not an attempt to legislate metaphysics. Rather, OI--EI is a constraint on theory language: do not build into a layer's ontology distinctions that the layer itself cannot, even in principle, stably witness. Spekkens' critique of many quantum ``solutions'' is that they preserve the standard formalism by paying with precisely such surplus structure.

\subsection{Formalization: empirical equivalence from a family of lenses}

Let $Z$ be a space of microdescriptions (``substrate states''). A layer is specified by a family of \emph{lenses}
\[
 f_i : Z \to X_i,\qquad i\in I,
\]
where $f_i(z)$ represents what is accessible to the layer via experiment or interface~$i$.
(For simplicity, one may take a common codomain~$X$ and write $f_i : Z\to X$; the construction below does not depend on this choice.)

Define an empirical equivalence relation:
\[
 z \sim z' \quad\Longleftrightarrow\quad \forall i\in I,\ f_i(z)=f_i(z').
\]
By construction, $z\sim z'$ means that the layer cannot distinguish $z$ from $z'$ using any admissible lens in the family. OI--EI then dictates that such states must be identified as the same object at that layer.

\subsection{The Leibniz quotient and its universal property}

The correct object space for the layer is therefore the quotient $Z/{\sim}$: an ``object'' at the layer is not a microstate, but an equivalence class of empirically indistinguishable microstates.
Moreover, each admissible lens must factor through this quotient, because by definition lens values are constant on equivalence classes.

\begin{theorem}[Leibniz quotient factorization]\label{thm:leibniz-quotient}
Let $Z$ be a type of microdescriptions and let $\{f_i : Z\to X\}_{i\in I}$ be a family of lenses. Define $z\sim z'$ iff $f_i(z)=f_i(z')$ for all $i$.
Then:
\begin{enumerate}
  \item The relation $\sim$ is an equivalence relation, hence determines a quotient map $q: Z\to Z/{\sim}$.
  \item Each lens $f_i$ factors through the quotient: there exists a unique $\bar f_i : Z/{\sim}\to X$ such that $\bar f_i\circ q = f_i$.
  \item (Universal property) If $g:Z\to Y$ is constant on $\sim$-classes, then there exists a unique $\bar g: Z/{\sim}\to Y$ with $\bar g\circ q = g$.
\end{enumerate}
\end{theorem}

The universal property itself is standard quotient mathematics; the substantive move is the methodological constraint that the lens family defines the layer's ontology. Theorem~\ref{thm:leibniz-quotient} is the precise mathematical statement of ``empirical indiscernibility implies ontological identity at the layer'': the layer's objects are exactly the equivalence classes, and every layer-expressible quantity factors uniquely through these classes. In the accompanying repository, this factorization is mechanized in Lean~4 (see \texttt{lean/Sbtq/Sbtq/LeibnizQuotient.lean}, specifically \texttt{quotient\_lift\_exists\_unique} and \texttt{liftLens\_comp\_q}; see Appendix~\ref{app:lean} for a consolidated table of mechanized statements).

\subsection{A finite example (mirroring the mechanization)}

Let $Z=\{\texttt{Bool}\}\times\{\texttt{Bool}\}$, and consider a layer with a single lens $f(z)=z_1$ that reads only the first bit. Then $(\texttt{true},\texttt{false})\sim(\texttt{true},\texttt{true})$ because the lens cannot access the second component. The Leibniz quotient has exactly two objects: the class whose first bit is \texttt{false} and the class whose first bit is \texttt{true}. Any layer-expressible predicate must therefore depend only on~$z_1$; treating the second bit as an object-level distinction would violate OI--EI for this layer.

This example provides the structural template that we will reuse for quantum theory: the relevant ``objects'' are determined by what can become a stable record at the layer, not by the full microdescription. The rest of the paper instantiates this quotient-and-factorization logic using packaging maps (closures) that stabilize record-level distinctions.
